\section*{Resolution, dynamics, and reconstruction}
A physical description often begins by separating a state into spatial cells, frequency bands, spin channels, tensor factors, or spectral sectors. Such a separation is useful only when the components can be reconstructed, their interactions can be described, and different resolutions can be compared. Those requirements lead to the calculus developed here.

An orthogonal projector selects a packet. Placing an operator between two projectors gives a transfer block. A skew-adjoint generator determines an antisymmetric current between packet weights. Grouping the labels produces a coarser description in which internal currents cancel. These statements can be understood in a small matrix example before any infinite-dimensional analysis is needed.

The same questions become more demanding when frames move, states become mixed, or the system exchanges information with an environment. Countable resolutions introduce convergence. Unbounded operators introduce domains. Each extension retains part of the finite algebra and adds a hypothesis that must remain visible where it is used.

The free Dirac construction that follows gives these ideas a concrete analytic realization. Fourier modes resolve a spinor field, finite matrices govern its components, and completion leads to a self-adjoint Hamiltonian. Gauge transport, constrained gravity, and arithmetic refinement then test which parts of the same language survive in different settings.

\begin{readerroadmap}[title={The organizing question}]
Which quantities can be reconstructed exactly, which balances survive a change of resolution, and which additional hypotheses are required for an infinite or unbounded system?
\end{readerroadmap}

The mathematical arguments use ordinary complex linear algebra, operator theory, and Hilbert-space analysis. Formal source correspondences and verification records are collected in the technical appendices. Numerical illustrations explain specified examples rather than replacing universal proofs.

\clearpage
\section*{Reading routes and conventions}
\addcontentsline{toc}{section}{Reading routes and conventions}
The first route follows the algebra. Begin with orthogonal resolution and block composition, then read the current identity, coarse aggregation, and moving frames. These sections provide the common notation for the rest of the volume. The worked examples keep the states and operators explicit.

The second route follows physical descriptions. Mixed states, tensor marginals, and Lindblad jumps show which information a packet weight retains and which information still resides in coherence. The positive-effect construction then explains how an overlapping resolution can be represented inside a larger sharp one. Classical probabilities arise only after the relevant normalization and closure assumptions have been supplied.

The third route follows analysis. Countable resolutions require an actual summation statement. An unbounded operator requires its domain. A strong limit acts on each fixed vector, while an operator-norm limit controls all unit vectors at once. The final chapters make these distinctions explicit and explain why the finite pointwise algebra $\C^N$ cannot support a nonzero complex-linear derivation obeying the ordinary Leibniz rule, even though matrix corners support nontrivial commutators.

\begin{longtable}{@{}L{.24\linewidth}L{.67\linewidth}@{}}
\toprule
Convention & Meaning\\\midrule\endhead
Inner product & $\inner{x}{y}$ is conjugate-linear in $x$ and linear in $y$.\\
Projector & $P_a=P_a^*=P_a^2$, with orthogonality between distinct packet labels.\\
Transfer block & $L_{ba}=P_bLP_a$ takes a component in packet $a$ to packet $b$.\\
State equation & $\dot\psi=K\psi$, with $K^*=-K$. The Hamiltonian convention is $K=-iH$ in units $\hbar=1$.\\
Packet weight & $p_a=\norm{P_a\psi}^2$. It is a probability for a normalized state and a complete orthogonal resolution. It is not an energy expectation by definition.\\
Current orientation & $J_{b\leftarrow a}$ is inflow into packet $b$ from packet $a$, so $\dot p_b=\sum_aJ_{b\leftarrow a}$ in the finite closed setting.\\
Mixed state & $\rho\geq0$, $\operatorname{tr}\rho=1$ in the finite-dimensional sections.\\
Coarse cell & A union of disjoint fine packet labels. Its projector is their sum, with the relevant convergence and domain hypotheses in the countable setting.\\
Unbounded action & An expression such as $H\psi$ includes the condition $\psi\in\Dom(H)$.\\
Time and units & $K$ has units of inverse time. We set $\hbar=1$ when identifying $K=-iH$. All plotted examples state their numerical parameters.\\
\bottomrule
\end{longtable}

Finite-dimensional statements appear with that hypothesis. A finite projector family can also act on an infinite-dimensional Hilbert space, provided the operators and differentiability assumptions in the statement make sense. These are different uses of the word finite.

The three companions serve different purposes. The monograph gives the mathematical argument. The visual atlas makes specific states, currents, and convergence examples visible. The Lean companion records exact declarations and evidence. A displayed ordinary proof is not silently promoted to a newly compiled theorem, and a sampled curve is not used to establish a universal limit.

\clearpage
